\section{行列群の共役類}
次の問題は, 意外だろう.
\begin{egprob} (H9数学I-3)
$A,B$を$n$次実正方行列とする. ある$P\in \mr{GL}_{n}(\mb{C})$が存在して$B=PAP^{-1}$であるとき, ある$Q\in \mr{GL}_{n}(\mb{R})$が存在して$B=QAQ^{-1}$であることを示せ.
\end{egprob}
\begin{egans}
$P$の実部を集めた行列$X$と, 虚部を集めた行列$Y$を考える. $X+iY = P$である. $t$を複素数として, $A+tB$を考える. $BP=PA$で, $A,B$が実であるから$BX=XA, BY=YA$である. よって$B(X+tY) = A(X+tY)$である. $\det{(X+tY)}$は$t$の多項式であるが, $t=i$のときに$\det{P} \neq 0$となるので, この多項式は零多項式とは異なる. したがって$t$を実数の範囲に制限して,$\det{(X+tY)}\neq 0$になるような$t$を見つけることができ, それを$Q$とおけばよい. \qed
\end{egans}




\subsection{(旧)$\mr{GL}_{2}(K)$の考察}
この節では$\mr{GL}_{2}(K)$の共役類について論じる. アイデアについては, \href{https://math.stackexchange.com/questions/1416055/classification-of-all-conjugacy-classes-of-gl-2-mathbbr-gl-2-mathbbq}{Stack}を参考にした.\par
Jordan標準形は, 代数閉体$K$に対して行列群の共役類を決定していた. たとえば$K=\mb{C}$で$\mr{GL}_{2}(\mb{C})$の共役類は, $\sumib{a&0\\ 0&b}$に代表されるもの, $\sumib{a&1\\ 0&a}$に代表されるものになる. この共役類上で最小多項式や固有多項式やトレースや行列式が一致する. \par
ところが代数閉体でない場合はやはり話が変わる. というのも, Jordan標準形にする行列$P$は固有ベクトルを並べて構成されたけれども, その固有ベクトルが$K$に値を持つかは不明で, $P$が考えている群に収まらないことがあるからだ.

\begin{eg} $\mr{GL}_{2}(\mb{R})$において$X=\sumib{1&2\\
-1&-1}$を考える. 固有多項式は$T^2+1$であり, 実根を持たない. $X-iI = \sumib{1-i & 2\\ -1 & -1-i}$の$\mb{C}$上の固有空間は$\sumib{1+i\\ -1}$で生成される1次元空間である. しかし, $\mb{R}$では固有ベクトルが\tf{存在しない.}
\end{eg}

では, 一般の体$K$上ではどのように分類するとよいか. Jordan標準形が考えられる部分はもちろんあるが, 問題はその外側である. 先ほどの例では, 固有値が$\mb{R}$に属さないことが問題であると思われた. 実際にここが大きな問題点である. なぜなら, 逆に, $X \in \mr{GL}_{2}(K)$の固有値$\lambda$が$K$に取れる場合は, $X-\lambda I$もまた$K$成分なので, それの核を求める(= 右辺0の連立方程式を解く)ことは$K$の中で行えるし, 固有値が相異なるかどうかに限らず, Jordan標準形にする$P$が$K$成分で取れることになる.  したがって, このときの固有値は再び共役類を決定する. \par
\begin{rem} \tf{GLの場合, 固有値が基礎体$K$に入らないのに$K$成分の固有ベクトルが取れるといったような"奇跡"みたいなことは起こらない.}  実際, $A\in \mr{GL}_{2}(K)$の固有値$\lambda\notin K$はノンゼロで, $v\in K^2\setminus{0}$に対し$Av = \lambda v$だとすると左辺は$K$成分だが, 右辺はそうではない. $v = \sumib{k_1 \\ k_2}$とでもおいてやれば, $\lambda$が$K$に入ってしまうことになるというのが分かるであろう.
\end{rem}
したがって, GLの場合には, $K$に固有ベクトルを持たないことと$K$内の固有値を持たないことは同値になる.

\lefba{
\begin{prop} \label{prop:when_eigen_doesnot_exist}
$K$は体とする.
$A\in \mr{GL}_{2}(K)$が$K$に固有ベクトルを持たないとする. このとき, $A$は次の行列に共役である.
\[\bmat{ 0 & -\det{A} \\ 1 & \tr{A}} \]
したがって, 同じ最小多項式を持つ$\mr{GL}_{2}(K)$の元同士は共役である.
\end{prop}
}
\prf ($P$の作り方は必読!) $v_1\in K^2$を0でないベクトルとし, $v_2 = Av_1$とする. \tf{$A$が固有ベクトルを$K$に持たないから, $v_1$と$v_2$は平行ではない.} したがって$v_1,v_2$は一次独立である. よって$P=\sumib{v_1 & v_2}$とおくと, $v_1,v_2$の一次独立性から$P\in \mr{GL}_{2}(K)$であって, $P$は標準基底$e_1,e_2$を$v_1,v_2$に送る変換であるから
\bealn{
P^{-1}AP &= P^{-1}\sumib{Av_1 & Av_2} = P^{-1}\sumib{v_2 & A^2v_1}\\
& P^{-1}\ao{= \sumib{v_2 & (\tr{A}A - (\det{A})I) v_1 }}\\
& = P^{-1}\sumib{v_2 & \tr{A} v_2 - (\det{A})v_1}\\
& = \sumib{e_2 & \tr{A} e_2 - (\det{A})e_1}\\
& = \bmat{0 & -\det{A} \\ 1 & \tr{A} }
}
\qed \\
ここで, \aka{$K$が標数2ではないとする.} このときは二次多項式に対して平方完成をすることができる.
一般に$K$係数の二次方程式$T^2 - pT + q = 0$が$K$に根を持つための$p,q$に関する必要十分条件は, $(T-p/2)^2 = \frac{p^2-4q}{2^2}$と変形すれば分かるように, $p^2-4q$が$K$の平方元($k\in K$を用いて$k^2$と表せる元)であること, である.
$T^2 - pT +q$を固有多項式に持つ行列$X\in \mr{GL}_{2}(K)$は, $q=0$では存在しないが($\because \det{X} \neq 0$), $q\neq 0$ならば, 先の命題の行列を模倣して
\[
\bmat{
0 & -q \\
1 & p
} \in \mr{GL}_{2}(K)
\]
というものが存在する. よって, $T^2- pT + q$で$p^2-4q$が平方元でないものに対し共役類が一つ割り当てられる. \par
最後に, $K$が\tf{標数2でなく}重複した固有値$\lambda$を持つ場合は, 固有値和$2\lambda$が$K$に属すので$\lambda\in K$となる. よって, それは$\sumib{\lambda & 0 \\ 0 & \lambda}, \sumib{\lambda & 1 \\ 0 & \lambda}\in \mr{GL}_{2}(K)$のどちらかに共役である. それぞれの場合で最小多項式(これは共役不変量！)は$T-\lambda$, $(T-\lambda)^2$である.
\par
以上をまとめる:

\lefba{
\begin{prop}  $K$は標数2でない体とする. $\mr{GL}_{2}(K)$の元$A$に対して次が成り立つ.
\ben
\item $A$の最小多項式が$K$の範囲で根$\lambda_1,\lambda_2$を持ち, 重根を持たないならば, $A$は$\sumib{\lambda_1 & 0 \\ 0 & \lambda_2}$に共役である(固有値の順番は問題にならない).
\item $A$が重複した固有値を持つ場合, 自動的にその固有値$\lambda$は$K$に属し, $A$は$\sumib{\lambda & 1\\ 0 & \lambda}$に共役である.
\item $A$の固有多項式が$T^2 - \tr{A} T + \det{A}$で, これが$K$の中に根を持たないならば,(自動的にその根は相異なり)$A$は $\sumib{0 & -\det{A} \\ 1 & \tr{A}}$に共役である.
\een
逆に, $(p,q)\in K^2$であって, $p^2 - 4q$が平方元でないかつ$q\neq 0$であるとき, 固有多項式が$T^2 - pT + q$であるような$\mr{GL}_{2}(K)$の元が存在する. たとえば, $\sumib{0 & -q \\ 1 & p}$がある.

\end{prop}
}
この命題を固有多項式から解釈すれば次のようになるであろう.
\begin{cor} (\tf{固有多項式から共役類を見つける方法})\\
$K$は標数2でない体, $A \in \mr{GL}_{2}(K)$, 固有多項式が$T^2 - pT + q$であるとする($q\neq 0$).
\ben
\item $p^2 - 4q$が$0$ではない平方元であるとき, $A\sim \sumib{\lambda_1 & 0 \\ 0 & \lambda_2} \sim \sumib{\lambda_2 & 0 \\ 0 & \lambda_1}$
\item $p^2 - 4q = 0$のとき, $A$の固有値は$p/2\in K$(重複度2)であり, $A\sim \sumib{p/2 & 0 \\ 0 & p/2}$ もしくは $A\sim \sumib{p/2 & 1\\ 0 & p/2}$である. この二つは最小多項式が$T-p/2$の何次であるかで分類できる.
\item $p^2 - 4q$が平方元でないとき, $A\sim \sumib{0 & -q \\ 1 & p}$.
\een
\end{cor}
\begin{egprob}
$\mb{R}$における平方元とは, まさしく0以上の実数である. $\mr{GL}_{2}(\mb{R})$の共役類をすべて決定せよ.
\end{egprob}
\begin{egans} $D = p^2-4q$とおく. $D>0$なら$T^2-pT + q$の相異なる根を$r_1,r_2$として$\sumib{r_1 & 0 \\ 0 & r_2}$の類$C_1$に属す. $D = 0$で最小多項式$T- p/2$なら$\sumib{p/2 & 0 \\ 0 & p/2}$の類$C_2$に属し, $(T-p/2)^2$なら$\sumib{p/2 & 1 \\ 0 & p/2}$の類$C_3$に属す. $D<0$なら$\sumib{0 & -q \\ 1 & p}$の類$C_{p,q}$に属す.
\end{egans}

$K=\mb{F}_{2}$の場合については最後に補足する.


過去問では有限体上で出題された.
\begin{egprob}

\end{egprob}

\begin{egans}
\chatgpthint{問題文が未記入です。}
\end{egans}
さて, 最後に$K=\mb{F}_{2^n}$の場合に$\mr{GL}_{2}(K)$の共役類について調べよう. 標数2なので, $T^2 - pT + q=0$が$K$に解を持つかはややこしい(平方元に関する特徴付けが効かなくなる). 固有多項式が相異なる2つの根を持つ場合か, 根を持たない場合については同様である. 重根を持つ場合を考えよう. $\lambda$が重根であるとすると, 固有多項式は$T^2 + \lambda^2$となる. $\lambda\notin K$でなくとも$\lambda^2 \in K$となることがある. もし$\lambda\in K$なら同様で, Jordanに直す$P$が$K$から取れる. $\lambda\notin K$の場合は固有ベクトルが$K$の中に取れないので, 同様の議論から$\sumib{0 & - \lambda^2 \\ 1 & 0}$に共役になる. したがって, 次のようにまとめられる.

\begin{prop} $K$が標数2の体であるとする. $A\in \mr{GL}_{2}(K)$とする.
\ben
\item $A$の最小多項式が$K$の範囲で相異なる根$\lambda_1,\lambda_2$を持つならば$A$は$\sumib{\lambda_1 & 0 \\ 0 & \lambda_2}$に共役である.
\item  $A$が\tf{$K$の範囲に}重複した固有値$\lambda\in K$を持つとする. もし最小多項式が$T-\lambda$なら$\sumib{\lambda & 0 \\ 0 & \lambda}$に共役で, $(T-\lambda)^2$なら$\sumib{\lambda & 1 \\ 0 & \lambda}$に共役である.
\item $A$が$K$の範囲に固有値を持たない場合は, $\sumib{0 & -\det{A} \\ 1 & \tr{A}}$に共役である. とくに, $\lambda$が重複した固有値で$\lambda\notin K$かつ$\lambda^2 \in K$を満たす場合, $\sumib{0 & -\lambda^2 \\ 1 & 0}$に共役である.
\een

\end{prop}



\subsection{$\mr{SL}_2(\mb{R})$の共役類}
$\mr{SL}_{2}(\mb{R})$の共役類を調べる. $A\in \mr{SL}_{2}(\mb{R})$とする. $A$が相異なる実根$r_1,r_2$を持つとき, $A$は$P\in \mr{GL}_{2}(\mb{R})$で対角化できる:
\[P^{-1} A P = \bmat{r_1 & 0 \\ 0 & r_2}\]
$P$がSLに入るか分からないが, $Q=\sumib{1&0 \\ 0 & -1/\det{P}}$とおくと$PQ\in \mr{SL}_{2}(\mb{R})$であり, $P^{-1} = \sumib{1 & 0 \\ 0 & -\det{P}}$である. よって,
\[(PQ)^{-1} A (PQ) = P^{-1}\bmat{r_1 & 0 \\ 0 & r_2}P = \bmat{r_1 & 0 \\ 0 & r_2}\]
となるので, $\mr{SL}_{2}(\mb{R})$においても$A\sim \sumib{r_1 & 0 \\ 0 & r_2}$であるとわかる. \par
$A$が重根$r$を持つとき, SLなので$r=\pm 1$である. この場合でも, もし$A$が対角化可能なら$\mr{SL}_{2}(\mb{R})$において$rI$に共役である. 対角化可能ではないとき, $P\in \mr{GL}_{2}(\mb{R})$が存在して
\[P^{-1}A P = \bmat{r & 1 \\ 0 & r}\]
となる. $P$に$\dfrac{1}{\sqrt{|\det{P}|}}$をかけて取り替えれば, $\det{P} = \pm 1$としてよい. 先と同じ$Q = \sumib{1 & 0 \\ 0 & \pm 1}$を取ると$PQ\in \mr{SL}_{2}(\mb{R})$である. $Q^{-1} = Q$なので,
\[
\bmat{
1 & 0 \\ 0 & \pm 1
} \bmat{r & 1 \\ 0 & r} \bmat{1 & 0 \\ 0 & \pm 1} = \bmat{r & \pm 1\\ 0 & r}
\]
だから, 結局$A$は右辺のどちらかに共役である.
{\small
\lefba{
\begin{prop} $r=\pm 1$のとき,
$P = \sumib{a&b \\ c&d}\in \mr{SL}_{2}(\mb{R})$によって$P\sumib{r & 1\\ 0 & r} = \sumib{r&-1 \\ 0 & r} P$となることはない. よって, これら二つの元で代表される共役類は相異なる.
\end{prop}
\prf
両辺計算すれば
\[\bmat{ar & a+br \\ cr & c+dr} = \bmat{ ar -c & br - d\\ cr & dr }\]
だから$c=0, a=-d$を得る. すると$ad-bc = -a^2 \neq 1$だから共役にならない.
\qed
}
}

最後に, $A$が虚数の固有値$z,\ovl{z}$を持つときを考える. この場合は命題\ref{prop:when_eigen_doesnot_exists}と同様に, $P$を上手く取って$P^{-1}AP=\sumib{0 & -\det{A} \\ 1 & \tr{A}} = \sumib{0 & -1 \\ 1 & \tr{A}}$である. この$P$がSLに属すとは限らないが, $\dfrac{1}{\sqrt{|\det{P}|}}$倍と取り替えて$\det{P}=\pm 1$としてよい. $\det{P}=1$なら$A$はこの行列と共役である. $\det{P}=-1$なら$Q = \sumib{1 & 0 \\ 0 & -1}$として, $PQ\in \mr{SL}_{2}(\mb{R})$であり,
\[(PQ)^{-1} A (PQ) = \bmat{0 & 1 \\ -1 & \tr{A}}\]
となるので$A$はこの右辺に共役である. こうして得られた二つの共役類が異なることが計算で示せる(\ref{linalg:prob:not_conjugate}). こうして, 次が得られた.
\lefba{
\begin{prop} ($\mr{SL}_{2}(\mb{R})$の共役類) $A\in \mr{SL}_{2}(\mb{R})$の属す共役類は次のようになる.
\ben
\item $A$が相異なる二つの実固有値$r_1,r_2$を持つ場合, $A\sim \sumib{r_1 & 0 \\ 0 & r_2}$.
\item $A$が重複する固有値$r$を持つ場合, $r=\pm 1$であって, もし$A$の最小多項式が$T-r$なら$A\sim rI$であり, $(T-r)^2$ならば$A\sim \sumib{r & 1 \\ 0 & r}$である.
\item $A$が虚数の固有値$z,\ovl{z}$を持つとき, $|z| = 1$, $-2 < \mr{Re}z < 2$であって, $P = \sumib{v_1 & Av_1}$に関して$\det{P}>0$であるなら$A\sim \sumib{0 & -1 \\ 1 & \tr{A}}$であり, $\det{P}<0$なら$A\sim \sumib{0 & 1 \\ -1 & \tr{A}}$である.
\een
したがって, $X$の共役類を$\set{X}$で表すと, 共役類は次ですべてである.
\[\parenb{\sumib{r & 0 \\ 0 & 1/r}}  \parenb{\sumib{1 & 1\\ 0 & 1}}, \parenb{\sumib{-1 & 1 \\ 0 & -1}}, \parenb{\sumib{0 & -1 \\ 1 & t}}, \parenb{\sumib{0 & 1\\ -1 & t}} )\]
ただし, $r\in \mb{R}$,
$-2 < t < 2$.
\end{prop}
}
検証してみよう.
\begin{egprob}
次の$\mr{SL}_{2}(\mb{R})$の元を上にリストした共役類のどれに属すかを調べよ(実際に共役であることを確かめること).
\[ \bmat{3&2\\ 4&3}, \bmat{1 & 0 \\ 1&1 }, \bmat{-1 & 0 \\ -1 & -1}, \bmat{1 & 2 \\ -1 & -1}, \bmat{-2 & 1\\ 1 & -1} \]

\end{egprob}




\subsection{演習問題}
\begin{prob} \label{linalg:prob:not_conjugate}
$A\in\mr{SL}_{2}(\mb{R})$が虚数の固有値$z,\ovl{z}$を持つとする. $\sumib{0 & 1 \\ -1 & \tr{A}}$と$\sumib{0 & -1 \\ 1 & \tr{A} }$は$\mr{SL}_{2}(\mb{R})$において共役ではないことを示せ.
\end{prob}
\begin{hint}
\chatgpthint{共役行列を未知成分で置き、行列式が1にならないことを示してください。}
\end{hint}
